%% ma: L10-B15-0049
%% lop: 10
%% chuong: 6
%% bai: 15
%% dang: 10.15.5
%% loai: TL
%% muc: TH
%% dap_an: "-"
%% nguon: "Ảnh giáo viên gửi 21/11/2026: Chuyên đề VI – Toán 10 – Chương VI (Hàm số – đồ thị và ứng dụng), tr. 199–203, Dạng 5 Câu 2"
%% kien_thuc: "Lớp 10 Bài 15 (hàm số đồng biến, nghịch biến)"
%% trang_thai: chinh-thuc
%% ghi_chu: "Dạng do giáo viên cung cấp (21/11/2026)."
\begin{ex}
Tìm tất cả các giá trị thực của tham số $m$ để hàm số $y=(2m+3)x+m+3$ nghịch biến trên $\mathbb R$.
\loigiai{$\dfrac{f(x_1)-f(x_2)}{x_1-x_2}=2m+3$. Hàm số nghịch biến $\Leftrightarrow2m+3<0\Leftrightarrow m<-\tfrac32$.}
\end{ex}
